% Fichier engendré par tools/generate-tex.py à partir de 01-nombres-calcul-algebrique.
% Ne pas éditer à la main : tout le texte provient des commentaires du fichier Lean.
\documentclass[11pt,a4paper]{article}

% Compile aussi bien avec pdflatex qu'avec xelatex ou tectonic.
\usepackage{iftex}
\ifPDFTeX
  \usepackage[T1]{fontenc}
  \usepackage[utf8]{inputenc}
\else
  \usepackage{fontspec}
\fi
\usepackage[french]{babel}
\usepackage{amsmath,amssymb,amsthm}
\usepackage[margin=2.5cm]{geometry}
\usepackage[hidelinks]{hyperref}

\theoremstyle{definition}
\newtheorem{definition}{Définition}
\theoremstyle{plain}
\newtheorem{theoreme}{Théorème}
\newtheorem{lemme}[theoreme]{Lemme}

% Un énoncé écrit mais non démontré : la preuve Lean est un `sorry`.
\newcommand{\admis}{\par\smallskip\noindent\textit{Admis : l'énoncé est écrit en Lean, sa démonstration reste à faire.}\par}

% Renvoi à la déclaration Lean, une ligne après la démonstration, aligné à droite.
\newcommand{\source}[2]{\par\smallskip\noindent\hfill{\footnotesize\href{#1}{\texttt{#2}}}\par\smallskip}

\title{Nombres, calcul algébrique}
\date{}

\begin{document}
\maketitle
% Les identifiants Lean en machine à écrire ne se coupent pas : on tolère des
% espacements plus lâches plutôt que des lignes qui débordent.
\sloppy

\section{NombresCalculAlgebrique}

\noindent
Lycée \textemdash{} section \og{} Nombres, calcul algébrique \fg{}.
Énoncés et démonstrations en français : voir NombresCalculAlgebrique.tex.

\subsection{Inclusions des ensembles de nombres}

\begin{theoreme}
Les entiers naturels sont des entiers relatifs, et l'inclusion est stricte : $-1$ n'est
pas un entier naturel.
\end{theoreme}

\begin{proof}
Tout entier naturel se lit comme entier relatif. Réciproquement, $-1$ n'est l'image
d'aucun entier naturel, les entiers naturels étant positifs.
\end{proof}
\source{https://github.com/Commutator-IO/learn-lean/blob/main/courses/02-lycee/01-nombres-calcul-algebrique/NombresCalculAlgebrique.lean\#L15}{NombresCalculAlgebrique.lean\#L15}

\begin{theoreme}
Les rationnels ne sont pas tous décimaux : $1/3$ n'a pas d'écriture décimale finie.
\end{theoreme}

\begin{proof}
Si $1/3 = a/10^n$, alors $10^n = 3a$, donc $3$ divise $10^n$. Or $3$ est premier et ne
divise pas $10$ : c'est impossible.
\end{proof}
\source{https://github.com/Commutator-IO/learn-lean/blob/main/courses/02-lycee/01-nombres-calcul-algebrique/NombresCalculAlgebrique.lean\#L21}{NombresCalculAlgebrique.lean\#L21}

\begin{theoreme}
Les réels ne sont pas tous rationnels.
\end{theoreme}

\begin{proof}
$\sqrt2$ en est un exemple, comme l'établit l'énoncé suivant.
\end{proof}
\source{https://github.com/Commutator-IO/learn-lean/blob/main/courses/02-lycee/01-nombres-calcul-algebrique/NombresCalculAlgebrique.lean\#L31}{NombresCalculAlgebrique.lean\#L31}

\subsection{\texttt{\(\sqrt{2}\)} est irrationnel}

\begin{theoreme}
$\sqrt2$ n'est pas rationnel : aucune fraction d'entiers n'a $2$ pour carré.
\end{theoreme}

\begin{proof}
Le raisonnement scolaire est celui de la parité : si $\sqrt2 = p/q$ avec la fraction
irréductible, alors $p^2 = 2q^2$, donc $p$ est pair, donc $q$ l'est aussi, ce qui contredit
l'irréductibilité.

Ce n'est pas celui que la bibliothèque formalise. Elle emploie une forme du théorème des
racines rationnelles : de $p^2 = 2q^2$ on tire que $q^2$ divise $p^2$ ; or, pour des
entiers, $q^n \mid p^n$ entraîne $q \mid p$, et l'irréductibilité force alors $q = 1$. Ainsi
$\sqrt2$ serait un entier, ce qu'il n'est pas puisque $1^2 = 1$ et $2^2 = 4$.

Les deux chemins concluent ; le second a l'avantage de valoir tel quel pour toute racine
$n$-ième d'un entier qui n'est pas une puissance $n$-ième exacte, sans qu'on ait à parler de
parité.
\end{proof}
\source{https://github.com/Commutator-IO/learn-lean/blob/main/courses/02-lycee/01-nombres-calcul-algebrique/NombresCalculAlgebrique.lean\#L36}{NombresCalculAlgebrique.lean\#L36}

\subsection{Valeur absolue et inégalité triangulaire}

\begin{theoreme}
La valeur absolue mesure une distance :
\[ |x - a| \le r \iff a - r \le x \le a + r . \]
\end{theoreme}

\begin{proof}
L'inégalité $|y| \le r$ équivaut à l'encadrement $-r \le y \le r$ ; en l'appliquant à
$y = x - a$ et en ajoutant $a$ aux trois membres, on obtient l'encadrement annoncé.
\end{proof}
\source{https://github.com/Commutator-IO/learn-lean/blob/main/courses/02-lycee/01-nombres-calcul-algebrique/NombresCalculAlgebrique.lean\#L42}{NombresCalculAlgebrique.lean\#L42}

\begin{theoreme}
Inégalité triangulaire : $|x + y| \le |x| + |y|$.
\end{theoreme}

\begin{proof}
Chacun des deux nombres est majoré par sa valeur absolue, et l'inégalité vaut aussi
pour les opposés ; en combinant les deux, la somme est encadrée par
$\pm(|x| + |y|)$.
\end{proof}
\source{https://github.com/Commutator-IO/learn-lean/blob/main/courses/02-lycee/01-nombres-calcul-algebrique/NombresCalculAlgebrique.lean\#L47}{NombresCalculAlgebrique.lean\#L47}

\subsection{Identités remarquables}

\begin{theoreme}
Identités remarquables :
\[ (a+b)^2 = a^2 + 2ab + b^2, \quad (a-b)^2 = a^2 - 2ab + b^2, \quad
   (a+b)(a-b) = a^2 - b^2 . \]
\end{theoreme}

\begin{proof}
Les trois se développent directement :
\[ (a+b)(a+b) = a^2 + ab + ba + b^2 = a^2 + 2ab + b^2 , \]
et de même pour les deux autres, les termes croisés s'annulant dans la troisième.
\end{proof}
\source{https://github.com/Commutator-IO/learn-lean/blob/main/courses/02-lycee/01-nombres-calcul-algebrique/NombresCalculAlgebrique.lean\#L52}{NombresCalculAlgebrique.lean\#L52}

\subsection{Produit nul et règle des signes}

\begin{theoreme}
Un produit est nul si et seulement si l'un des facteurs est nul.
\end{theoreme}

\begin{proof}
Si l'un des facteurs est nul, le produit l'est. Réciproquement, si $a \ne 0$, diviser
$ab = 0$ par $a$ donne $b = 0$.
\end{proof}
\source{https://github.com/Commutator-IO/learn-lean/blob/main/courses/02-lycee/01-nombres-calcul-algebrique/NombresCalculAlgebrique.lean\#L61}{NombresCalculAlgebrique.lean\#L61}

\begin{theoreme}
Un quotient de dénominateur positif est positif si et seulement si son numérateur
l'est.
\end{theoreme}

\begin{proof}
Multiplier par le dénominateur, strictement positif, conserve le sens de
l'inégalité.
\end{proof}
\source{https://github.com/Commutator-IO/learn-lean/blob/main/courses/02-lycee/01-nombres-calcul-algebrique/NombresCalculAlgebrique.lean\#L64}{NombresCalculAlgebrique.lean\#L64}

\subsection{Signe de \texttt{ax + b}}

\begin{theoreme}
Signe d'une expression du premier degré : pour $a > 0$,
\[ ax + b > 0 \iff x > -\frac{b}{a}, \qquad ax + b < 0 \iff x < -\frac{b}{a} . \]
\end{theoreme}

\begin{proof}
On isole $x$ en divisant par $a$, strictement positif, ce qui ne change pas le sens des
inégalités. La racine $-b/a$ sépare donc la droite en deux demi-droites de signes
contraires.
\end{proof}
\source{https://github.com/Commutator-IO/learn-lean/blob/main/courses/02-lycee/01-nombres-calcul-algebrique/NombresCalculAlgebrique.lean\#L70}{NombresCalculAlgebrique.lean\#L70}

\subsection{Puissances et racines}

\begin{theoreme}
La racine d'un produit de nombres positifs est le produit des racines :
$\sqrt{ab} = \sqrt a \sqrt b$.
\end{theoreme}

\begin{proof}
Les deux membres sont positifs et ont le même carré, à savoir $ab$.
\end{proof}
\source{https://github.com/Commutator-IO/learn-lean/blob/main/courses/02-lycee/01-nombres-calcul-algebrique/NombresCalculAlgebrique.lean\#L81}{NombresCalculAlgebrique.lean\#L81}

\begin{theoreme}
Puissance d'une puissance : $(a^n)^m = a^{nm}$.
\end{theoreme}

\begin{proof}
$(a^n)^m$ est le produit de $m$ facteurs valant chacun $a^n$, soit $nm$ facteurs
$a$.
\end{proof}
\source{https://github.com/Commutator-IO/learn-lean/blob/main/courses/02-lycee/01-nombres-calcul-algebrique/NombresCalculAlgebrique.lean\#L85}{NombresCalculAlgebrique.lean\#L85}

\subsection{Comparaison de \texttt{x}, \texttt{x²} et \texttt{\(\sqrt{x}\)}}

\begin{theoreme}
Sur $[0 ; 1]$, élever au carré diminue et prendre la racine augmente :
$x^2 \le x \le \sqrt x$.
\end{theoreme}

\begin{proof}
Pour la première inégalité, $x^2 \le x$ s'écrit $x(x - 1) \le 0$, vrai car $x \ge 0$ et
$x \le 1$.

Pour la seconde, posons $s = \sqrt x$, de sorte que $x = s^2$ et $0 \le s \le 1$. Alors
$x = s^2 \le s = \sqrt x$, par la première inégalité appliquée à $s$.
\end{proof}
\source{https://github.com/Commutator-IO/learn-lean/blob/main/courses/02-lycee/01-nombres-calcul-algebrique/NombresCalculAlgebrique.lean\#L91}{NombresCalculAlgebrique.lean\#L91}

\begin{theoreme}
Au-delà de $1$, c'est l'inverse : $\sqrt x \le x \le x^2$.
\end{theoreme}

\begin{proof}
Même argument avec les inégalités renversées : $x(x-1) \ge 0$ donne $x \le x^2$, et
$s = \sqrt x \ge 1$ donne $s \le s^2 = x$.
\end{proof}
\source{https://github.com/Commutator-IO/learn-lean/blob/main/courses/02-lycee/01-nombres-calcul-algebrique/NombresCalculAlgebrique.lean\#L103}{NombresCalculAlgebrique.lean\#L103}

\subsection{Sommes}

\begin{theoreme}
Somme des $n$ premiers entiers :
\[ 2\,(0 + 1 + \cdots + n) = n(n+1) . \]
\end{theoreme}

\begin{proof}
C'est le calcul de Gauss : on écrit la somme deux fois, une fois à l'endroit et une fois
à l'envers, et on ajoute terme à terme. Chacune des $n+1$ colonnes vaut $n$, d'où le
double de la somme.

L'énoncé est écrit sans division pour rester dans les entiers : $n(n+1)$ est toujours
pair, mais l'écriture $n(n+1)/2$ obligerait à parler de division euclidienne.
\end{proof}
\source{https://github.com/Commutator-IO/learn-lean/blob/main/courses/02-lycee/01-nombres-calcul-algebrique/NombresCalculAlgebrique.lean\#L116}{NombresCalculAlgebrique.lean\#L116}

\begin{theoreme}
Somme géométrique : pour $q \ne 1$,
\[ 1 + q + \cdots + q^n = \frac{1 - q^{n+1}}{1 - q} . \]
\end{theoreme}

\begin{proof}
En multipliant la somme $S$ par $1 - q$, tous les termes se télescopent sauf le premier
et le dernier :
\[ (1-q)S = (1 + q + \cdots + q^n) - (q + \cdots + q^{n+1}) = 1 - q^{n+1} , \]
puis on divise par $1 - q$, non nul puisque $q \ne 1$.
\end{proof}
\source{https://github.com/Commutator-IO/learn-lean/blob/main/courses/02-lycee/01-nombres-calcul-algebrique/NombresCalculAlgebrique.lean\#L122}{NombresCalculAlgebrique.lean\#L122}

\subsection{Coefficients binomiaux et binôme de Newton}

\begin{theoreme}
Relation de Pascal : $\binom{n}{k} + \binom{n}{k+1} = \binom{n+1}{k+1}$.
\end{theoreme}

\begin{proof}
C'est la règle de construction du triangle de Pascal : pour choisir $k+1$ éléments parmi
$n+1$, ou bien on prend le dernier — restent $k$ éléments à choisir parmi $n$ — ou bien on
ne le prend pas, et il faut en choisir $k+1$ parmi $n$.
\end{proof}
\source{https://github.com/Commutator-IO/learn-lean/blob/main/courses/02-lycee/01-nombres-calcul-algebrique/NombresCalculAlgebrique.lean\#L132}{NombresCalculAlgebrique.lean\#L132}

\begin{theoreme}
Formule du binôme de Newton :
\[ (a+b)^n = \sum_{k=0}^{n} \binom{n}{k} a^k b^{n-k} . \]
\end{theoreme}

\begin{proof}
En développant le produit des $n$ facteurs $(a+b)$, chaque terme est obtenu en
choisissant $a$ dans $k$ facteurs et $b$ dans les $n-k$ autres ; il y a $\binom{n}{k}$
façons de faire ce choix.
\end{proof}
\source{https://github.com/Commutator-IO/learn-lean/blob/main/courses/02-lycee/01-nombres-calcul-algebrique/NombresCalculAlgebrique.lean\#L137}{NombresCalculAlgebrique.lean\#L137}

\subsection{Raisonnement par récurrence}

\begin{theoreme}
Principe de récurrence : une propriété vraie au rang zéro et héréditaire est vraie à
tout rang.
\end{theoreme}

\begin{proof}
C'est le principe même de construction des entiers naturels : tout entier s'obtient à
partir de zéro par des passages au successeur, et l'hypothèse d'hérédité permet de suivre
cette construction. En Lean, la démonstration est le récurseur des entiers, c'est-à-dire
la définition du type.
\end{proof}
\source{https://github.com/Commutator-IO/learn-lean/blob/main/courses/02-lycee/01-nombres-calcul-algebrique/NombresCalculAlgebrique.lean\#L146}{NombresCalculAlgebrique.lean\#L146}

\subsection{Inégalité de Bernoulli}

\begin{theoreme}
Inégalité de Bernoulli : $(1 + a)^n \ge 1 + na$ pour $a \ge -1$.
\end{theoreme}

\begin{proof}
Par récurrence sur $n$. Au rang $0$, les deux membres valent $1$.

Si l'inégalité vaut au rang $n$, on multiplie par $1 + a \ge 0$ :
\[ (1+a)^{n+1} \ge (1 + na)(1 + a) = 1 + (n+1)a + na^2 \ge 1 + (n+1)a , \]
le terme $na^2$ étant positif. C'est l'exemple canonique de démonstration par récurrence
du programme, et il sert à établir que $q^n$ tend vers l'infini pour $q > 1$.
\end{proof}
\source{https://github.com/Commutator-IO/learn-lean/blob/main/courses/02-lycee/01-nombres-calcul-algebrique/NombresCalculAlgebrique.lean\#L153}{NombresCalculAlgebrique.lean\#L153}

\subsection{Énoncés admis}

\noindent
Ce que le programme demande et que ce document ne démontre pas. Les énoncés sont écrits en Lean, pour que le manque se compte ; leur démonstration est admise.

\begin{theoreme}
Le développement décimal d'un rationnel positif est périodique à partir d'un certain
rang : il existe une période $p > 0$ et un rang $n_0$ tels que, pour tout $n \ge n_0$, le
chiffre de rang $n + p$ égale celui de rang $n$.
\end{theoreme}
\admis
\source{https://github.com/Commutator-IO/learn-lean/blob/main/courses/02-lycee/01-nombres-calcul-algebrique/NombresCalculAlgebrique.lean\#L174}{NombresCalculAlgebrique.lean\#L174}

\vfill
\noindent\rule{0.35\linewidth}{0.4pt}\par
\noindent{\footnotesize Leonardo de Moura et Sebastian Ullrich,
\emph{The Lean~4 Theorem Prover and Programming Language},
\emph{Automated Deduction — CADE~28}, Springer, 2021, p.~625--635,
\href{https://doi.org/10.1007/978-3-030-79876-5_37}{doi:10.1007/978-3-030-79876-5\_37}.\par}

\end{document}
